\section{Etat de l'art}

\subsection{Le genre Visual Novel}

Le \gls{vn} est un genre de jeu vidéo interactif qui combine narration, visuels, musiques et parfois doublage vocal. Le genre est historiquement ancré dans la culture japonaise, avec des centaines de titres produits chaque année, et a progressivement gagné en popularité en Occident, notamment grâce à la montée en popularité des adaptations animées dans les années 2000.

Sur le plan technique, un \gls{vn} repose sur plusieurs mécaniques fondamentales: une narration textuelle qui progresse par clics, des images statiques représentant les personnages et les décors, des choix de dialogues, et une emphase marquée sur le design sonore et la bande originale.

Camingue, Carstensdottir et Melcer ont relevé que neuf des trente définitions académiques du genre mentionnent la présence de narration à embranchements. Ces structures peuvent aboutir à des expériences très variées: \textit{Steins;Gate 0} (2015) propose six fins différentes, tandis que \textit{428:Shibuya Scramble} (2008) peut mener à pas moins de 85 conclusions distinctes. À l'inverse, 18\% des titres recensés ne comportent aucun embranchement narratif. \cite{camingue_what_2021}

Du point de vue du marché, la période 2019-2024 a connu une croissance de la popularité du genre, porté par l'essor des plateformes de distributions comme Steam et les boutiques d'applications mobiles. Sur la plateforme Steam, nous comptons 16 985 jeux du genre sur Steam en Suisse\footnote{Nombre récupéré au 31 mars 2026}. Le marché mondial du \gls{vn} était estimé à 139 millions de dollars en 2023 et devrait atteindre 620 millions d'ici 2032, avec un taux de croissance annuel moyen de 8,5\%. \cite{noauthor_visual_nodate}



\subsection{Moteur de Visual Novels existants}

Les moteurs de jeu-videos sont des outils permettant de faciliter la création de ces derniers.

Nous pouvons parler de RenPy, un moteur de \gls{vn} utilisés par des milliers de personnes, offrant une interface graphique facile à prendre en main. RenPy permet le développement à destination des ordinateurs comme des téléphones portables.



\subsection{Approches no-code / low-code}

L'objectif de ce type de programmation est de faciliter l'entrée dans la programmation. Ou de permettre à l'aide d'outils graphiques de s'abstraire de l'écriture de code.

L'approche no-code (traduction: pas de code) et low-code (traduction: peux de code) sont très similaires, mais reste différentes. L'une demande toujours quelques connaissances en code ou du moins cela pourrait être utile, tandis que l'autre ne demande aucune connaissances et peut être utilisé directement par tous.
